\begin{table}[htbp]
\centering
\footnotesize
\caption{分季截距（各附自助法 95\% 置信区间）与季间异质性，$S=4$ 与 $S=3$ 分别报告。}
\label{tab:s5}
\begin{tabular}{llll}
\toprule
流行季 & $R_A$，规则窗口（$S=4$） & $G$，规则窗口（$S=4$） & $G$，日历窗口（$S=3$） \\
\midrule
2022-23 & 0.0137 [0.0066, 0.0213] & 0.0718 [0.0368, 0.1181] & 0.0711 [0.0347, 0.1124] \\
2023-24 & 0.0062 [0.0025, 0.0144] & 0.0305 [0.0119, 0.0700] & 0.0277 [0.0037, 0.0776] \\
2024-25 & 0.0127 [0.0057, 0.0241] & 0.0647 [0.0315, 0.1181] & 0.1352 [0.0602, 0.1857] \\
2025-26 & 0.0225 [0.0087, 0.0352] & 0.1383 [0.0426, 0.2285] & -- \\
$a_{\mathrm{RE}}$ [95\% CI] & 0.0125 [0.0022, 0.0227] & 0.0619 [0.0066, 0.1171] & 0.0735 [$-$0.0558, 0.2028] \\
$Q$ ($p$) & 5.71 (0.126) & 6.58 (0.087) & 8.92 (0.012) \\
$10^3\tau^2$ [Q-profile] & 0.019 [0.000, 0.599] & 0.620 [0.000, 26.756] & 2.077 [0.127, 114.859] \\
$I^2$ & 0.47 & 0.54 & 0.78 \\
\bottomrule
\end{tabular}
\par\smallskip
\begin{minipage}{0.95\linewidth}\scriptsize 分季截距以该季自身均值计算；95\% 双向自助区间，$B=1000$。随机效应采用 DerSimonian--Laird 估计及 $t_{S-1}$ 区间；$\tau^2$ 给出 Q-profile 95\% 区间。由 $S$ 季估计的季间方差，其相对标准误至少为 $\sqrt{2/(S-1)}$：$S=3$ 时为 1.00，$S=4$ 时为 0.82。\end{minipage}
\end{table}
